
\documentclass[11pt]{article}
\usepackage[margin=1in]{geometry}
\usepackage[T1]{fontenc}
\usepackage[utf8]{inputenc}
\usepackage{lmodern}
\usepackage{xcolor}
\usepackage{hyperref}
\usepackage{enumitem}
\usepackage{titlesec}
\usepackage{fancyhdr}
\usepackage{listings}
\usepackage{parskip}

\definecolor{navy}{RGB}{31,78,121}
\definecolor{lightgray}{RGB}{245,247,250}
\definecolor{codegray}{RGB}{70,70,70}

\hypersetup{
  colorlinks=true,
  linkcolor=navy,
  urlcolor=navy
}

\pagestyle{fancy}
\fancyhf{}
\lhead{FRED Dashboard Template Guide}
\rhead{\thepage}

\titleformat{\section}{\large\bfseries\color{navy}}{\thesection}{0.5em}{}
\titleformat{\subsection}{\normalsize\bfseries\color{navy}}{\thesubsection}{0.5em}{}

\lstset{
  basicstyle=\ttfamily\small\color{codegray},
  backgroundcolor=\color{lightgray},
  frame=single,
  breaklines=true,
  columns=fullflexible
}

\title{\textbf{Beginner-Friendly FRED Dashboard Template Guide}}
\author{}
\date{}

\begin{document}
\maketitle

\section{Purpose of this template}
This template follows the same overall format as the example dashboard project you uploaded.
It is designed so that a beginner can duplicate the structure, fill in a new dataset topic,
and still keep a professional modular dashboard architecture.

\section{Folder structure}
\begin{lstlisting}
project/
├── app.py
├── config.py
├── narrative_blocks.py
├── requirements.txt
├── .env.example
├── README.md
├── .streamlit/
│   └── config.toml
└── src/
    ├── data_loader.py
    ├── analysis.py
    └── visuals.py
\end{lstlisting}

\section{What each file means}

\subsection{app.py}
This is the main Streamlit dashboard file. It controls:
\begin{itemize}[leftmargin=1.5em]
  \item page layout
  \item sidebar controls
  \item top metric cards
  \item charts
  \item model output
  \item dataset download button
\end{itemize}
Most beginners should avoid changing app.py unless they want to add a brand new section.

\subsection{config.py}
This is the main customization file. For most projects, this is the only file that needs to be edited.
Replace every placeholder that says \texttt{FILL IN}. The most important parts are:
\begin{itemize}[leftmargin=1.5em]
  \item project title
  \item research question
  \item FRED series dictionary
  \item target variable
  \item model features
  \item chart groups
  \item start date
  \item target frequency
\end{itemize}

\subsection{src/data_loader.py}
This file downloads data from FRED and aligns mixed-frequency time series.
If your charts look blank, this is one of the first places to inspect.

\subsection{src/analysis.py}
This file handles summary statistics and the baseline linear regression model.
This project uses linear regression because the target is numeric.

\subsection{src/visuals.py}
This file creates the charts. Keeping chart code here makes app.py much cleaner.

\subsection{narrative\_blocks.py}
This file contains optional text helper functions. It is useful if you want reusable story text.

\subsection{requirements.txt}
This file lists the Python packages that must be installed.

\subsection{.env.example}
This shows the format for your private FRED API key. Copy it to a file named \texttt{.env}.

\section{How to customize the dashboard}

\subsection{Step 1: choose your FRED topic}
Pick one main variable and several supporting variables. Example topics:
\begin{itemize}[leftmargin=1.5em]
  \item unemployment
  \item inflation
  \item housing
  \item interest rates
  \item GDP and recession
  \item federal deficit
\end{itemize}

\subsection{Step 2: edit config.py}
In \texttt{config.py}, replace:
\begin{itemize}[leftmargin=1.5em]
  \item \texttt{PROJECT\_TITLE}
  \item \texttt{RESEARCH\_QUESTION}
  \item \texttt{SERIES}
  \item \texttt{TARGET\_VARIABLE}
  \item \texttt{MODEL\_FEATURES}
  \item \texttt{CHART\_GROUPS}
  \item \texttt{START\_DATE}
  \item \texttt{TARGET\_FREQUENCY}
\end{itemize}

\subsection{Step 3: create your .env file}
Create a file named \texttt{.env} in the project root and add:
\begin{lstlisting}
FRED_API_KEY=your_actual_key_here
\end{lstlisting}

\subsection{Step 4: install dependencies}
Mac or Linux:
\begin{lstlisting}
python3 -m venv .venv
source .venv/bin/activate
pip install -r requirements.txt
\end{lstlisting}

Windows PowerShell:
\begin{lstlisting}
python -m venv .venv
.venv\Scripts\Activate.ps1
pip install -r requirements.txt
\end{lstlisting}

\subsection{Step 5: run the dashboard}
\begin{lstlisting}
streamlit run app.py
\end{lstlisting}

\section{How the files work together}
The dashboard runs in this order:
\begin{enumerate}[leftmargin=1.5em]
  \item app.py reads settings from config.py
  \item app.py calls src/data\_loader.py to fetch and merge data
  \item app.py renames raw FRED IDs into nice display names
  \item app.py calls src/visuals.py to draw charts
  \item app.py calls src/analysis.py to compute metrics and fit the model
  \item the dashboard renders in Streamlit
\end{enumerate}

\section{Where the machine learning model is}
The baseline model is inside \texttt{src/analysis.py}. It uses linear regression:
\begin{lstlisting}
model = LinearRegression()
model.fit(X, y)
\end{lstlisting}
This is not logistic regression. Logistic regression is used for classification.
Linear regression is appropriate here because the target variable is numeric.

\section{Common beginner mistakes}
\begin{itemize}[leftmargin=1.5em]
  \item using fake FRED IDs instead of real ones
  \item forgetting to create the .env file
  \item making TARGET\_VARIABLE different from the keys in SERIES
  \item choosing a start date that is too early for one of the series
  \item mixing annual and monthly data without understanding target frequency
  \item leaving debug tables turned on in the final presentation
\end{itemize}

\section{Final pre-submission checklist}
\begin{itemize}[leftmargin=1.5em]
  \item all FILL IN placeholders replaced
  \item .env file created
  \item streamlit app runs without crashing
  \item charts display real lines
  \item regression table appears
  \item debug tables turned off
  \item findings paragraph updated with your actual conclusions
\end{itemize}

\end{document}
